\section*{Capítulo \ref{ch:b2:heart} --- El corazón y el ciclo cardíaco}

\begin{solution}{exo:b2:heart:1}
\omterm{def:b2:blood-circulation:vessels}{Vena} cava $\to$ \omterm{def:b2:heart:anatomy}{aurícula} derecha $\to$ válvula tricúspide $\to$ \omterm{def:b2:heart:anatomy}{ventrículo}
derecho $\to$ válvula pulmonar $\to$ \omterm{def:b2:blood-circulation:vessels}{arteria} pulmonar $\to$ pulmones
$\to$ \omterm{def:b2:blood-circulation:vessels}{venas} pulmonares $\to$ \omterm{def:b2:heart:anatomy}{aurícula} izquierda $\to$ válvula mitral $\to$
\omterm{def:b2:heart:anatomy}{ventrículo} izquierdo $\to$ válvula aórtica $\to$ aorta.
\end{solution}

\begin{solution}{exo:b2:heart:2}
Llenado (diástole): válvulas AV abiertas, semilunares cerradas.
Contracción isovolumétrica: las cuatro cerradas. Eyección: semilunares
abiertas, AV cerradas. Relajación isovolumétrica: todas cerradas. Primer
ruido: el cierre de las válvulas AV al comienzo de la sístole; segundo:
el cierre de las válvulas semilunares al final.
\end{solution}

\begin{solution}{exo:b2:heart:3}
\omterm{prop:b2:heart:conduction}{Nódulo sinoauricular} ($t = 0$; la onda P al despolarizarse las
\omterm{def:b2:heart:anatomy}{aurículas}); \omterm{prop:b2:heart:conduction}{nódulo auriculoventricular} (al que se llega hacia los
\qty{80}{ms} y que retiene el impulso \qty{100}{ms}: el segmento PR); haz
de His y \omterm{prop:b2:heart:conduction}{fibras de Purkinje} (\qtyrange{180}{220}{ms}; el QRS al
despolarizarse los \omterm{def:b2:heart:anatomy}{ventrículos}); más tarde, la onda T al repolarizarse.
\end{solution}

\begin{solution}{exo:b2:heart:4}
Dentro de su intervalo de funcionamiento, el \omterm{def:b2:heart:anatomy}{ventrículo} expulsa más cuanto
más se ha llenado. Si el \omterm{def:b2:heart:anatomy}{ventrículo} derecho bombea más, el izquierdo
recibe más unos latidos después y, por la misma ley, bombea más: todo
desequilibrio se corrige solo, sin ninguna señal.
\end{solution}

\begin{solution}{exo:b2:heart:5}
\omterm{prop:b2:heart:cycle}{Volumen sistólico} de \qty{80}{mL}; \omterm{prop:b2:heart:cycle}{fracción de eyección},
$80/130 = \qty{62}{\%}$; gasto, $80\times 70 = \qty{5.6}{L/min}$. Insuficiente:
\qty{60}{mL}, $60/180 = \qty{33}{\%}$, \qty{4.2}{L/min} --- el \omterm{def:b2:heart:anatomy}{ventrículo}
dilatado mantiene casi el gasto trabajando con un gran volumen, pero solo
expulsa un tercio de lo que contiene.
\end{solution}

\begin{solution}{exo:b2:heart:6}
$W = 100\times 133\times 7\times 10^{-5} = \qty{0.93}{J}$; potencia,
$0.93\times 1.2 = \qty{1.1}{W}$. A \qty{150}{mmHg}: \qty{1.4}{J},
\qty{1.7}{W}. Los \qty{0.56}{W} mecánicos adicionales son \qty{3.7}{W}
metabólicos, \qty{0.19}{mL} de oxígeno por segundo: \qty{11}{mL/min} más.
\end{solution}

\begin{solution}{exo:b2:heart:7}
$60/0.5 = 120$ latidos por minuto; $60/1.5 = 40$. Tercero: bloqueo
cardíaco completo --- las \omterm{def:b2:heart:anatomy}{aurículas} laten a 100 bajo el mando del nódulo,
los \omterm{def:b2:heart:anatomy}{ventrículos} a 33 bajo un marcapasos propio, y el nódulo AV no conduce
nada.
\end{solution}

\begin{solution}{exo:b2:heart:8}
Deriva: $20/25 = \qty{0.8}{s}$, más \qty{0.15}{s}: \qty{0.95}{s}, 63
latidos por minuto. Acetilcolina: $25/18 = \qty{1.39}{s}$, más $0.15$:
\qty{1.54}{s}, 39 por minuto. Noradrenalina: $20/40 = \qty{0.5}{s}$,
más $0.15$: \qty{0.65}{s}, 92 por minuto.
\end{solution}

\begin{solution}{exo:b2:heart:9}
El retraso permite que las \omterm{def:b2:heart:anatomy}{aurículas} terminen de vaciarse en los
\omterm{def:b2:heart:anatomy}{ventrículos} antes de que estos se contraigan; sin él, \omterm{def:b2:heart:anatomy}{aurículas} y
\omterm{def:b2:heart:anatomy}{ventrículos} se contraerían a la vez, las válvulas AV se cerrarían sobre
unos \omterm{def:b2:heart:anatomy}{ventrículos} medio llenos y el \omterm{prop:b2:heart:cycle}{volumen sistólico} perdería la
aportación auricular. Un intervalo PR largo solo retrasa la sístole
ventricular y no cuesta nada hasta que el retraso se alarga tanto que se
pierden latidos o la conducción falla por completo.
\end{solution}

\begin{solution}{exo:b2:heart:10}
Reposo: $5000/45 = \qty{111}{mL}$; ejercicio: $\num{30000}/180 =
\qty{167}{mL}$. La ley de Starling convierte el mayor retorno venoso del
ejercicio en un mayor \omterm{prop:b2:heart:cycle}{volumen sistólico}; los nervios simpáticos añaden
frecuencia y contractilidad (un volumen telesistólico menor). Por encima
de unos 200, la diástole es demasiado corta para llenar el \omterm{def:b2:heart:anatomy}{ventrículo}, y el
\omterm{prop:b2:heart:cycle}{volumen sistólico} cae más deprisa de lo que sube la frecuencia.
\end{solution}

\begin{solution}{exo:b2:heart:11}
$v = Q/A$: $250/3 = \qty{83}{cm/s}$ y $250/1 = \qty{250}{cm/s}$.
$\Delta P = \tfrac12\times 1060\times v^{2}$: \qty{365}{Pa}
(\qty{2.7}{mmHg}) y \qty{3300}{Pa} (\qty{25}{mmHg}). El \omterm{def:b2:heart:anatomy}{ventrículo} genera
la presión adicional engrosando su pared; el \omterm{def:b2:heart:anatomy}{ventrículo} grueso y rígido
acaba superando su riego coronario, se llena mal y fracasa.
\end{solution}

\begin{solution}{exo:b2:heart:12}
El automatismo da un latido sin ninguna entrada; la ley de Starling ajusta
el gasto al retorno y equilibra los dos \omterm{def:b2:heart:anatomy}{ventrículos} sin ninguna entrada;
los nervios y la adrenalina solo fijan la frecuencia y la contractilidad
en torno a ese núcleo autorregulado. Un \omterm{def:b2:heart:anatomy}{corazón} trasplantado late (a unos
100, sin tono vagal), ajusta su \omterm{prop:b2:heart:cycle}{volumen sistólico} al llenado y puede
aumentar su gasto durante el ejercicio mediante el mecanismo de Starling y
la adrenalina circulante --- pero despacio, y menos que un \omterm{def:b2:heart:anatomy}{corazón}
inervado.
\end{solution}

\begin{solution}{pb:b2:heart:1}
\textbf{1.} \qty{70}{mL}; $70/130 = \qty{54}{\%}$; $70\times 70 =
\qty{4.9}{L/min}$.
\textbf{2.} $60/70 = \qty{0.86}{s}$; diástole, $0.86 - 0.3 =
\qty{0.56}{s}$.
\textbf{3.} $70/0.3 = \qty{233}{mL/s}$; $233/3 = \qty{78}{cm/s}$.
\textbf{4.} $70/1500 = \qty{0.047}{s}$, unos \qty{50}{ms}.
\textbf{5.} $100\times 133\times 7\times 10^{-5} = \qty{0.93}{J}$;
$0.93\times 70/60 = \qty{1.1}{W}$.
\textbf{6.} Ambos \omterm{def:b2:heart:anatomy}{ventrículos}: $1.1\times 1.2 = \qty{1.3}{W}$; con un
\qty{15}{\%}, \qty{8.7}{W} metabólicos; $8.7/20 = \qty{0.43}{mL}$ de
oxígeno por segundo, \qty{26}{mL/min}.
\textbf{7.} Extrayendo \qty{0.14}{mL} por mililitro de \omterm{def:b2:blood-circulation:blood}{sangre}:
$26/0.14 = \qty{190}{mL/min}$.
\textbf{8.} Intervalo de \qty{0.86}{s}, tiempo de deriva de \qty{0.71}{s}:
$20/0.71 = \qty{28}{mV/s}$.
\textbf{9.} Intervalo de \qty{0.6}{s}, deriva de \qty{0.45}{s}:
\qty{44}{mV/s}.
\textbf{10.} Intervalo de \qty{0.33}{s}, deriva de \qty{0.18}{s}:
\qty{110}{mV/s}.
\textbf{11.} R--R de \qty{0.86}{s}; P: despolarización auricular; QRS:
despolarización ventricular; T: repolarización ventricular.
\textbf{12.} El retraso permite que las \omterm{def:b2:heart:anatomy}{aurículas} se vacíen en los
\omterm{def:b2:heart:anatomy}{ventrículos} antes de que estos se contraigan; a 180, el ciclo completo
dura \qty{0.33}{s}, de modo que un retraso invariable de \qty{0.16}{s}
consumiría la mitad --- los nervios simpáticos aceleran también la
conducción del nódulo.
\textbf{13.} Bloqueo cardíaco completo: los \omterm{def:b2:heart:anatomy}{ventrículos} a 40 por minuto,
con el ritmo marcado por el haz o las \omterm{prop:b2:heart:conduction}{fibras de Purkinje}; las \omterm{def:b2:heart:anatomy}{aurículas}, a
75 bajo el mando del \omterm{prop:b2:heart:conduction}{nódulo sinoauricular}.
\textbf{14.} \qty{120}{mL}; $120/150 = \qty{80}{\%}$; $120\times 180 =
\qty{21.6}{L/min}$.
\textbf{15.} $0.33 - 0.3 = \qty{0.03}{s}$: casi no hay tiempo para llenarse;
más deprisa aún, y el \omterm{prop:b2:heart:cycle}{volumen sistólico} se desploma.
\textbf{16.} $W = 120\times 133\times 1.2\times 10^{-4} = \qty{1.9}{J}$,
$\times 3$ por segundo: \qty{5.7}{W}; ambos \omterm{def:b2:heart:anatomy}{ventrículos}, \qty{6.9}{W};
metabólicos, \qty{46}{W}; oxígeno, \qty{2.3}{mL/s}, \qty{140}{mL/min}.
\textbf{17.} $140/0.14 = \qty{1000}{mL/min}$, cinco veces el flujo en
reposo, que debe llegar en una diástole reducida a la décima parte del
ciclo: las \omterm{def:b2:blood-circulation:vessels}{arteriolas} coronarias deben dilatarse al máximo.
\textbf{18.} Frecuencia de unos 100 sin tono vagal, volumen telesistólico
sin cambios en \qty{60}{mL}: \omterm{prop:b2:heart:cycle}{volumen sistólico} de $150 - 60 = \qty{90}{mL}$,
gasto de \qty{9}{L/min}.
\textbf{19.} Solo la frecuencia ($70 \to 180$ con \qty{70}{mL}): $+7.7$;
el llenado (\qty{20}{mL} más a 180): $+3.6$; la contractilidad
(\qty{30}{mL} menos de volumen residual a 180): $+5.4$ --- de 4,9 a
\qty{21.6}{L/min}: los tres incrementos dan cuenta exactamente de la
subida.
\textbf{20.} $233/1 = \qty{233}{cm/s}$, \qty{2.3}{m/s}.
\textbf{21.} $\tfrac12\times 1060\times 2.33^{2} = \qty{2900}{Pa}$,
\qty{22}{mmHg}.
\textbf{22.} $120/0.3 = \qty{400}{mL/s}$, \qty{4}{m/s}; $\tfrac12\times
1060\times 16 = \qty{8500}{Pa}$, \qty{64}{mmHg}.
\textbf{23.} Máxima de $120 + 22 = \qty{142}{mmHg}$ en reposo, y de unos
$140 +
64 = \qty{204}{mmHg}$ durante el ejercicio; la presión media de eyección en
reposo sube de 100 a 122: trabajo por latido $\times 1.22$.
\textbf{24.} Una pared más gruesa genera más fuerza y reduce la tensión
por fibra; pero la masa que hay que irrigar crece mientras que el flujo
coronario, comprimido por la pared gruesa en sístole y con una diástole
más corta, no sigue el ritmo, y una pared gruesa se relaja mal y se llena
menos: siguen la isquemia y la insuficiencia.
\textbf{25.} \qty{4.9}{L/min} en reposo, \qty{21.6}{L/min} durante el
ejercicio; \qty{0.93}{J} por latido; gradiente de \qty{22}{mmHg} en reposo
y de \qty{64}{mmHg} durante el ejercicio.
\end{solution}
